\pagestyle{empty}
\tight
\begin{tblr}{width=\textwidth,colspec={|Q[c,m,1.9cm]|X[l,m]|},hlines,vlines,hspan=minimal,row{1}={bg=headblue},rowsep=1pt,colsep=3pt}
\SetCell{c,m}\hfil\logo\hfil & \textbf{POLITEKNIK NEGERI MALANG}\par\textbf{JURUSAN TEKNOLOGI INFORMASI}\par\textbf{PROGRAM STUDI: D4 TEKNIK INFORMATIKA} \\
\SetCell[c=2]{c} \textbf{RENCANA PEMBELAJARAN SEMESTER (RPS)} & \\
\end{tblr}
\begin{longtblr}{width=\textwidth,colspec={|Q[l,m,1.9cm]|X[0.85,l,m]|X[1.45,l,m]|X[0.75,l,m]|X[1.15,l,m]|},hlines,vlines,hspan=minimal,rowsep=1pt,colsep=2pt,row{1,3}={bg=headgray,font=\bfseries}}
MATA KULIAH & KODE & BOBOT (SKS)/jam & SEMESTER & TGL. PENYUSUNAN \\
Pembelajaran Mesin & RTI255004 & 3 SKS/6 jam & 5 & 1 Juli 2025 \\
OTORISASI & Dosen Pengembang RPS & Koordinator RMK & \SetCell[c=2]{l} Ka PRODI &  \\
 & Mamluatul Hani'ah, S.Kom., M.Kom. & Mamluatul Hani'ah, S.Kom., M.Kom. & \SetCell[c=2]{l}{\vspace{8mm}\par Dr. Ely Setyo Astuti, S.T., M.T.} &  \\
\SetCell[r=9]{l} Capaian Pembelajaran (CP) & \SetCell[c=4]{l,bg=headgray,font=\bfseries} Capaian Pembelajaran Lulusan Program Studi (CPL-Prodi) &  &  &  \\
 & CPL07 & \SetCell[c=3]{l} Mampu mengembangkan sistem komputasi cerdas dan melakukan analisis data berbasis kecerdasan buatan untuk pengambilan keputusan. &  &  \\
 & CPL10 & \SetCell[c=3]{l} Mampu menganalisis persoalan kompleks dengan wawasan multidisiplin untuk menghasilkan solusi di bidang informatika dan ilmu komputer. &  &  \\
 & \SetCell[c=4]{l,bg=headgray,font=\bfseries} Capaian Pembelajaran mata kuliah (CPL-MK) &  &  &  \\
 & CPMK0706 & \SetCell[c=3]{l} Mahasiswa mampu mengembangkan sistem ML dan melakukan analisis data berbasis AI untuk menghasilkan keputusan yang bermakna. &  &  \\
 & CPMK1005 & \SetCell[c=3]{l} Mahasiswa mampu mengevaluasi solusi sistem cerdas berdasarkan kinerja, keandalan, dan interpretabilitas model. &  &  \\
 & \SetCell[c=4]{l,bg=headgray,font=\bfseries} Kemampuan akhir tiap tahapan belajar (Sub-CPMK) &  &  &  \\
 & SCPMK0706-04001 & \SetCell[c=3]{l} Mahasiswa mampu mengimplementasikan algoritma ML supervised dan unsupervised untuk menyelesaikan masalah klasifikasi, regresi, dan clustering. &  &  \\
 & SCPMK1005-04002 & \SetCell[c=3]{l} Mahasiswa mampu mengevaluasi kinerja model ML menggunakan metrik yang tepat dan menerapkan teknik optimasi model. &  &  \\
\SetCell[r=2]{l} Penilaian & \SetCell[c=4]{l} *Nilai CPMK didapat dari agregasi nilai SUB-CPMK yang dibebankan &  &  &  \\
 & \SetCell[c=4]{l}{\tiny\begin{tblr}{width=\linewidth,colspec={|Q[c,m,0.1\linewidth]|Q[c,m,0.16\linewidth]|X[l,m]|X[l,m]|X[l,m]|X[l,m]|X[l,m]|X[l,m]|Q[c,m,0.08\linewidth]|},hlines,vlines,hspan=minimal,rowsep=1pt,colsep=0.7pt,row{1,2}={bg=headgray,font=\bfseries}}
Kode CPMK & Kode SCPMK & \SetCell[c=6]{c} Bobot per Bentuk Penilaian & & & & & & Total Bobot \\
 & & Tugas & Kuis 1 & UTS & Kuis 2 & PBL & UAS & \\
CPMK0706 & SCPMK0706-04001 & 16\%: implementasi algoritma ML & 5\%: supervised learning & 10\%: supervised learning dan evaluasi & 0\% & 12\%: end-to-end ML project & 10\%: evaluasi ML project & 53\% \\
CPMK1005 & SCPMK1005-04002 & 12\%: evaluasi dan optimasi model & 0\% & 5\%: supervised learning dan evaluasi & 5\%: unsupervised dan optimasi & 13\%: end-to-end ML project & 12\%: evaluasi ML project & 47\% \\
\SetCell[c=8]{r}\textbf{Total Per Penilaian} & & & & & & & & \textbf{100\%} \\
\end{tblr}} &  &  &  \\
Diskripsi Singkat Mata Kuliah & \SetCell[c=4]{l} Pembelajaran Mesin membangun kemampuan mengimplementasikan, mengevaluasi, dan mengoptimalkan algoritma ML untuk klasifikasi, regresi, clustering, dan deep learning dasar menggunakan scikit-learn dan Python dalam proyek end-to-end. &  &  &  \\
Bahan Kajian / Materi Pembelajaran & \SetCell[c=4]{l}\begin{enumerate}\item Supervised learning: regresi, decision tree, random forest, SVM, k-NN\item Model evaluation: cross-validation, confusion matrix, metrik, regularisasi\item Unsupervised learning: k-Means, DBSCAN, dan dimensionality reduction (PCA)\item Ensemble methods, deep learning dasar, dan interpretabilitas model\item End-to-end ML project: pipeline, optimasi, dan presentasi\end{enumerate} &  &  &  \\
\SetCell[r=2]{l} Pustaka & \SetCell[c=4]{l} \textbf{Utama:}\begin{enumerate}\item Géron, A. 2022. Hands-On Machine Learning (3rd ed.). O'Reilly.\item Raschka, S. 2022. Machine Learning with PyTorch and Scikit-Learn. Packt.\item James, G. 2021. An Introduction to Statistical Learning. Springer.\end{enumerate} &  &  &  \\
 & \SetCell[c=4]{l} \textbf{Pendukung:} Materi LMS, dokumentasi resmi scikit-learn dan PyTorch, dan jobsheet praktikum. &  &  &  \\
Media Pembelajaran & \SetCell[c=2]{l} \textbf{Software:} Python, scikit-learn, Jupyter Notebook, LMS. &  & \SetCell[c=2]{l} \textbf{Hardware:} Komputer/laptop, proyektor. &  \\
Nama Dosen Pengampu & \SetCell[c=4]{l} Tim Pengajar Program Studi D4 TI &  &  &  \\
Matakuliah Syarat & \SetCell[c=4]{l} RTI254008 Statistik Komputasi &  &  &  \\
\end{longtblr}
\begin{longtblr}{width=\textwidth,colspec={|Q[c,m,0.06\textwidth]|X[1.35,l,m]|X[1.08,l,m]|X[1.16,l,m]|X[0.7,l,m]|X[1.24,l,m]|X[1.06,l,m]|X[0.98,l,m]|Q[c,m,0.05\textwidth]|},hlines,vlines,hspan=minimal,rowhead=2,row{1,2}={bg=headgreen,font=\bfseries},rowsep=1pt,colsep=1.5pt}
Mgg.\par Ke & Kemampuan Akhir Yang Direncanakan (Sub-CP-MK) & Bahan kajian (Materi Pembelajaran) & Bentuk dan Metode Pembelajaran & Estimasi Waktu & Pengalaman Belajar Mahasiswa & Kriteria \& Bentuk Penilaian & Indikator Penilaian & Bobot Penilaian (\%) \\
(1) & (2) & (3) & (4) & (5) & (6) & (7) & (8) & (9) \\
1 & SCPMK0706-04001 & Pengantar ML: paradigma, pipeline, dan ekosistem Python. & Modalitas: Blended Learning. Bentuk: Luring/Daring. Metode: Case Method, praktik ML, studi kasus, dan peer review. & 1 x 3 x 50' tatap muka; 1 x 3 x 50' tugas/praktik mandiri. & Mahasiswa mengerjakan aktivitas terarah untuk pengantar ML: paradigma, pipeline, dan ekosistem Python dan menunjukkan evidence hasil belajar. & Bentuk penilaian: Pengantar ML. Mengacu rubrik RTM. & Ketepatan konsep; kualitas implementasi model; validasi dan dokumentasi. & 2 \\
2 & SCPMK0706-04001 & Exploratory Data Analysis dan preprocessing. & Modalitas: Blended Learning. Bentuk: Luring/Daring. Metode: Case Method, praktik ML, studi kasus, dan peer review. & 1 x 3 x 50' tatap muka; 1 x 3 x 50' tugas/praktik mandiri. & Mahasiswa mengerjakan aktivitas terarah untuk exploratory data analysis dan preprocessing dan menunjukkan evidence hasil belajar. & Bentuk penilaian: Exploratory Data Analysis dan Preprocessing. Mengacu rubrik RTM. & Ketepatan konsep; kualitas implementasi model; validasi dan dokumentasi. & 2 \\
3 & SCPMK0706-04001 & Linear regression dan logistic regression. & Modalitas: Blended Learning. Bentuk: Luring/Daring. Metode: Case Method, praktik ML, studi kasus, dan peer review. & 1 x 3 x 50' tatap muka; 1 x 3 x 50' tugas/praktik mandiri. & Mahasiswa mengerjakan aktivitas terarah untuk linear regression dan logistic regression dan menunjukkan evidence hasil belajar. & Bentuk penilaian: Linear Regression dan Logistic Regression. Mengacu rubrik RTM. & Ketepatan konsep; kualitas implementasi model; validasi dan dokumentasi. & 2 \\
4 & SCPMK0706-04001 & Decision tree dan random forest. & Modalitas: Blended Learning. Bentuk: Luring/Daring. Metode: Case Method, praktik ML, studi kasus, dan peer review. & 1 x 3 x 50' tatap muka; 1 x 3 x 50' tugas/praktik mandiri. & Mahasiswa mengerjakan aktivitas terarah untuk decision tree dan random forest dan menunjukkan evidence hasil belajar. & Bentuk penilaian: Decision Tree dan Random Forest. Mengacu rubrik RTM. & Ketepatan konsep; kualitas implementasi model; validasi dan dokumentasi. & 2 \\
5 & SCPMK0706-04001 & Kuis 1: supervised learning. & Modalitas: Blended Learning. Bentuk: Luring/Daring. Metode: Case Method, praktik ML, studi kasus, dan peer review. & 1 x 3 x 50' tatap muka; 1 x 3 x 50' tugas/praktik mandiri. & Mahasiswa mengerjakan aktivitas terarah untuk kuis 1: supervised learning dan menunjukkan evidence hasil belajar. & Bentuk penilaian: Kuis 1: Supervised Learning. Mengacu rubrik RTM. & Ketepatan konsep; kualitas implementasi model; validasi dan dokumentasi. & 5 \\
6 & SCPMK0706-04001 & SVM dan k-NN classifier. & Modalitas: Blended Learning. Bentuk: Luring/Daring. Metode: Case Method, praktik ML, studi kasus, dan peer review. & 1 x 3 x 50' tatap muka; 1 x 3 x 50' tugas/praktik mandiri. & Mahasiswa mengerjakan aktivitas terarah untuk SVM dan k-NN classifier dan menunjukkan evidence hasil belajar. & Bentuk penilaian: SVM dan k-NN Classifier. Mengacu rubrik RTM. & Ketepatan konsep; kualitas implementasi model; validasi dan dokumentasi. & 2.5 \\
7 & SCPMK1005-04002 & Model evaluation: cross-validation, confusion matrix. & Modalitas: Blended Learning. Bentuk: Luring/Daring. Metode: Case Method, praktik ML, studi kasus, dan peer review. & 1 x 3 x 50' tatap muka; 1 x 3 x 50' tugas/praktik mandiri. & Mahasiswa mengerjakan aktivitas terarah untuk model evaluation: cross-validation, confusion matrix dan menunjukkan evidence hasil belajar. & Bentuk penilaian: Model Evaluation. Mengacu rubrik RTM. & Ketepatan konsep; kualitas implementasi model; validasi dan dokumentasi. & 2.5 \\
8 & SCPMK0706-04001, SCPMK1005-04002 & UTS: supervised learning dan evaluasi. & Modalitas: Blended Learning. Bentuk: Luring/Daring. Metode: Case Method, praktik ML, studi kasus, dan peer review. & 1 x 3 x 50' tatap muka; 1 x 3 x 50' tugas/praktik mandiri. & Mahasiswa mengerjakan aktivitas terarah untuk UTS: supervised learning dan evaluasi dan menunjukkan evidence hasil belajar. & Bentuk penilaian: UTS: Supervised Learning dan Evaluasi. Mengacu rubrik RTM. & Ketepatan konsep; kualitas implementasi model; validasi dan dokumentasi. & 15 \\
9 & SCPMK0706-04001 & Unsupervised learning: k-Means dan DBSCAN. & Modalitas: Blended Learning. Bentuk: Luring/Daring. Metode: Case Method, praktik ML, studi kasus, dan peer review. & 1 x 3 x 50' tatap muka; 1 x 3 x 50' tugas/praktik mandiri. & Mahasiswa mengerjakan aktivitas terarah untuk unsupervised learning: k-Means dan DBSCAN dan menunjukkan evidence hasil belajar. & Bentuk penilaian: Unsupervised Learning: k-Means dan DBSCAN. Mengacu rubrik RTM. & Ketepatan konsep; kualitas implementasi model; validasi dan dokumentasi. & 2.5 \\
10 & SCPMK0706-04001 & Dimensionality reduction: PCA. & Modalitas: Blended Learning. Bentuk: Luring/Daring. Metode: Case Method, praktik ML, studi kasus, dan peer review. & 1 x 3 x 50' tatap muka; 1 x 3 x 50' tugas/praktik mandiri. & Mahasiswa mengerjakan aktivitas terarah untuk dimensionality reduction: PCA dan menunjukkan evidence hasil belajar. & Bentuk penilaian: Dimensionality Reduction: PCA. Mengacu rubrik RTM. & Ketepatan konsep; kualitas implementasi model; validasi dan dokumentasi. & 2.5 \\
11 & SCPMK0706-04001 & Ensemble methods: Gradient Boosting, XGBoost. & Modalitas: Blended Learning. Bentuk: Luring/Daring. Metode: Case Method, praktik ML, studi kasus, dan peer review. & 1 x 3 x 50' tatap muka; 1 x 3 x 50' tugas/praktik mandiri. & Mahasiswa mengerjakan aktivitas terarah untuk ensemble methods: Gradient Boosting, XGBoost dan menunjukkan evidence hasil belajar. & Bentuk penilaian: Ensemble Methods: Gradient Boosting, XGBoost. Mengacu rubrik RTM. & Ketepatan konsep; kualitas implementasi model; validasi dan dokumentasi. & 2.5 \\
12 & SCPMK1005-04002 & Hyperparameter tuning dan regularisasi. & Modalitas: Blended Learning. Bentuk: Luring/Daring. Metode: Case Method, praktik ML, studi kasus, dan peer review. & 1 x 3 x 50' tatap muka; 1 x 3 x 50' tugas/praktik mandiri. & Mahasiswa mengerjakan aktivitas terarah untuk hyperparameter tuning dan regularisasi dan menunjukkan evidence hasil belajar. & Bentuk penilaian: Hyperparameter Tuning dan Regularisasi. Mengacu rubrik RTM. & Ketepatan konsep; kualitas implementasi model; validasi dan dokumentasi. & 2.5 \\
13 & SCPMK1005-04002 & Kuis 2: unsupervised learning dan optimasi. & Modalitas: Blended Learning. Bentuk: Luring/Daring. Metode: Case Method, praktik ML, studi kasus, dan peer review. & 1 x 3 x 50' tatap muka; 1 x 3 x 50' tugas/praktik mandiri. & Mahasiswa mengerjakan aktivitas terarah untuk kuis 2: unsupervised learning dan optimasi dan menunjukkan evidence hasil belajar. & Bentuk penilaian: Kuis 2: Unsupervised Learning dan Optimasi. Mengacu rubrik RTM. & Ketepatan konsep; kualitas implementasi model; validasi dan dokumentasi. & 5 \\
14 & SCPMK0706-04001 & Pengantar deep learning: ANN dan backpropagation. & Modalitas: Blended Learning. Bentuk: Luring/Daring. Metode: Case Method, praktik ML, studi kasus, dan peer review. & 1 x 3 x 50' tatap muka; 1 x 3 x 50' tugas/praktik mandiri. & Mahasiswa mengerjakan aktivitas terarah untuk pengantar deep learning: ANN dan backpropagation dan menunjukkan evidence hasil belajar. & Bentuk penilaian: Pengantar Deep Learning: ANN dan Backpropagation. Mengacu rubrik RTM. & Ketepatan konsep; kualitas implementasi model; validasi dan dokumentasi. & 2.5 \\
15 & SCPMK1005-04002 & Interpretabilitas model: SHAP, LIME. & Modalitas: Blended Learning. Bentuk: Luring/Daring. Metode: Case Method, praktik ML, studi kasus, dan peer review. & 1 x 3 x 50' tatap muka; 1 x 3 x 50' tugas/praktik mandiri. & Mahasiswa mengerjakan aktivitas terarah untuk interpretabilitas model: SHAP, LIME dan menunjukkan evidence hasil belajar. & Bentuk penilaian: Interpretabilitas Model: SHAP, LIME. Mengacu rubrik RTM. & Ketepatan konsep; kualitas implementasi model; validasi dan dokumentasi. & 2.5 \\
16 & SCPMK0706-04001, SCPMK1005-04002 & PBL: end-to-end ML project. & Modalitas: Blended Learning. Bentuk: Luring/Daring. Metode: Case Method, praktik ML, studi kasus, dan peer review. & 1 x 3 x 50' tatap muka; 1 x 3 x 50' tugas/praktik mandiri. & Mahasiswa mengerjakan aktivitas terarah untuk PBL: end-to-end ML project dan menunjukkan evidence hasil belajar. & Bentuk penilaian: PBL: End-to-End ML Project. Mengacu rubrik RTM. & Ketepatan konsep; kualitas implementasi model; validasi dan dokumentasi. & 25 \\
17 & SCPMK1005-04002 & UAS: presentasi dan evaluasi ML project. & Modalitas: Blended Learning. Bentuk: Luring/Daring. Metode: Case Method, praktik ML, studi kasus, dan peer review. & 1 x 3 x 50' tatap muka; 1 x 3 x 50' tugas/praktik mandiri. & Mahasiswa mengerjakan aktivitas terarah untuk UAS: presentasi dan evaluasi ML project dan menunjukkan evidence hasil belajar. & Bentuk penilaian: UAS: Presentasi dan Evaluasi ML Project. Mengacu rubrik RTM. & Ketepatan konsep; kualitas implementasi model; validasi dan dokumentasi. & 22 \\
\end{longtblr}
